% Generated draft: requires mathematical and source review.
\documentclass[11pt]{article}
\usepackage[margin=25mm]{geometry}
\usepackage{fontspec}
\setmainfont{Latin Modern Roman}
\usepackage{amsmath,amssymb,mathtools,unicode-math}
\setmathfont{Latin Modern Math}
\usepackage{enumitem,needspace,xurl,hyperref}
\setlength{\parindent}{0pt}
\setlength{\parskip}{5pt}
\setlength{\emergencystretch}{3em}
\newcommand{\N}{\mathbb{N}}
\newcommand{\Z}{\mathbb{Z}}
\newcommand{\Q}{\mathbb{Q}}
\newcommand{\R}{\mathbb{R}}
\newcommand{\C}{\mathbb{C}}
\newcommand{\F}{\mathbb{F}}
\newcommand{\A}{\mathbb{A}}
\newcommand{\PP}{\mathbb{P}}
\newcommand{\E}{\mathbb{E}}
\newcommand{\eps}{\varepsilon}
\newcommand{\dd}{\,\mathrm{d}}
\begin{document}
\section*{UnsolvedMath ID: 20001875}
\textbf{Source problem code:} catalogue entry
\subsection{Definitions and mathematical statement}
Paternain's source question, in the reported surrounding setting of magnetic flows on closed surfaces, asks whether every energy level has at least one closed magnetic geodesic. Use a closed Riemannian surface \((M,g)\), a closed magnetic two-form \(\sigma\), and kinetic energy \(H(x,p)=|p|_g^2/2\) on \(T^*M\) with the magnetic twist of its canonical symplectic form. Fix the sign convention by the equivalent tangent equation
\[
 \nabla_t\dot\gamma=Y(\dot\gamma),\qquad
 g(Yv,w)=\sigma(v,w).
\]
The question concerns periodic trajectories on regular positive levels \(H=E>0\); the zero level consists of stationary points and is not the relevant regular level.

The supplied counterexample takes a closed oriented hyperbolic surface of constant curvature \(-\kappa^2\), \(\kappa>0\), and \(\sigma=b\,\dd A_g\), \(b\ne0\). This form is not exact, since its integral over \(M\) is nonzero. At speed \(s=\sqrt{2E}\), magnetic trajectories have signed geodesic curvature \(b/s\). The aid's precise trichotomy is:
\begin{itemize}
\item If \(0<s<|b|/\kappa\), lifted trajectories are hyperbolic circles. Every trajectory is periodic and contractible, with prime physical-time period
\[
 T_{\mathrm{sub}}(s)=\frac{2\pi}{\sqrt{b^2-\kappa^2s^2}}.
\]
\item If \(s=|b|/\kappa\), lifted trajectories are horocycles. No trajectory is periodic on the closed quotient: a period would require a parabolic element in its torsion-free cocompact surface lattice, which has none. Thus the regular positive level
\[
 E_c=\frac{b^2}{2\kappa^2}
\]
contains no closed magnetic geodesic, refuting the universal assertion as printed.
\item If \(s>|b|/\kappa\), every nontrivial free homotopy class has a closed magnetic geodesic. For a primitive class whose hyperbolic geodesic representative has length \(\ell\), its period is
\[
 T_{\mathrm{sup}}(s;\ell)=
 \frac{\kappa\ell}{\sqrt{\kappa^2s^2-b^2}}.
\]
Nonprimitive classes give the corresponding iterates. Geometrically these trajectories are hypercycles at the appropriate signed distance from the geodesic axis.
\end{itemize}
For fixed nonzero \(b\), and fixed \(\ell\) on the supercritical side, both period expressions diverge with inverse-square-root order as the critical speed is approached. This counterexample does not contradict an existence theorem assuming exactness of \(\sigma\) or additional contact-type hypotheses.

The original discussion attributes an affirmative exact-form result to Contreras--Macarini--Paternain [Co-Ma-Pa2004], partly based on [Taimanov1992], proposes applying Hofer--Wysocki--Zehnder [Ho-Wy-Ze1993] after establishing the appropriate contact type, and cites the survey [Ginzburg1996]. Their full hypotheses are not included. Nonexactness obstructs a globally defined magnetic Lagrangian using a primitive one-form on the surface; it does not mean that no local or covering-space variational formulation can exist. The contact-type proposal is a conditional route, not proof that every level is contact type.

The extracted record then splices in the next section on path and loop spaces. With \(\Lambda M\) the loop space and \(\Omega(p,q)\) the fixed-endpoint path space, it discusses length sublevels \(\Lambda^TM\) and \(\Omega^T(p,q)\), energy critical points as closed or connecting geodesics, and finite-complex/finite-Betti-number assertions attributed to [Milnor1963]. These are background for a separate subject, not added hypotheses or conclusions of the magnetic question. The extraction's assertion that loops of bounded length themselves ``correspond to critical points'' is imprecise: only geodesics are energy critical. Compactness, regularity and sublevel conventions for the finite-complex claim are not supplied, so no unrestricted theorem about arbitrary Riemannian manifolds and all levels is asserted here.
\subsection{Short English statement}
A closed hyperbolic surface with a constant nonexact magnetic field has a regular energy level whose trajectories are horocycles and none are periodic. Below this level all trajectories are periodic circles; above it nontrivial free homotopy classes have closed trajectories. The original universal existence question therefore fails in the supplied setting.
\subsection{Sources}
\begin{itemize}
\item Permanent UnsolvedMath ID \(20001875\), exact public code catalogue entry; local registry \texttt{aim-geometry-notes.json}, \texttt{notes[212]}, canonical tag question.
\item AIM workshop ``Geodesics,'' Geometry, section 2.4 ``Of magnetic flows on closed surfaces,'' Question 2.4.1, attributed to Paternain: \url{https://aimath.org/pastworkshops/geodesicsproblems.pdf}. The section context is reported by the supplied verification note; no PDF was accessed for this draft.
\item Original explicit reference labels: Contreras, Macarini and Paternain [Co-Ma-Pa2004]; [Taimanov1992]; Hofer, Wysocki and Zehnder [Ho-Wy-Ze1993]; survey [Ginzburg1996]; and [Milnor1963] in the spliced path-space background. Full titles and publication details are absent, and the separate original literature field is empty.
\item Research attempt 1, ``A critical horocycle energy level without closed magnetic geodesics,'' labelled counterexample, refuted at the run, proved and low novelty confidence; registry solved and SOLVED-IN-LITERATURE. The statement remains reconstructed\_unverified. The source-verification field reports the surrounding closed-surface Hamiltonian setting rather than supplying a new full citation.
\end{itemize}
\end{document}
